% !TeX program = xelatex
% 本文件可正常編譯，請編譯兩次；刻意錯誤另見 warnings-zh-TW.tex。
\documentclass[UTF8,fontset=none,11pt]{ctexrep}
% Use an installed Traditional Chinese font; no font files are bundled.
\newcommand{\setexampleCJK}[1]{
  \setCJKmainfont[AutoFakeSlant=0.2]{#1}
  \setCJKsansfont{#1}
  \setCJKmonofont{#1}
}
\IfFontExistsTF{Noto Serif CJK TC}{\setexampleCJK{Noto Serif CJK TC}}{
  \IfFontExistsTF{Noto Serif TC}{\setexampleCJK{Noto Serif TC}}{
    \IfFontExistsTF{PingFang TC}{\setexampleCJK{PingFang TC}}{
      \PackageError{sample-font}{Install Noto Serif CJK TC or Noto Serif TC}
        {This sample requires an installed Traditional Chinese font.}
    }
  }
}
\usepackage{amsmath}
\usepackage{xcolor}
\usepackage{listings}
\usepackage[hidelinks]{hyperref}
\setcounter{secnumdepth}{5}
\setcounter{tocdepth}{3}
\lstdefinestyle{sample}{
  basicstyle=\ttfamily\small,
  numbers=left,stepnumber=1,numberfirstline=true,numberblanklines=false,
  numbersep=8pt,frame=single,breaklines=true
}
\lstset{style=sample}
\title{LaTeX Structure Highlight：繁體中文範例}
\author{示範文件}
\date{}

\begin{document}
\maketitle
\tableofcontents
\chapter{更容易閱讀的文件}
\label{chap:overview}
\section{實驗設計}
\label{sec:design}
每一層標題在編輯器中都有不同顏色。編譯兩次後，可以在標題與引用上方看到實際編號。
這些標示不會改變 PDF 的外觀。

\subsection{資料準備}
\label{sec:data}
檔名可以寫成 \texttt{data\_file.csv}，跳脫後的底線是合法的。
行內字面文字 \verb|raw_name| 也合法。

\subsubsection{輸入檢查}
\label{sec:validation}
數學下標 $x_i$ 與 \(y_j\) 不應被當成錯誤底線。
本範例以 \texttt{secnumdepth=5} 開啟至 subparagraph 的編號。

\paragraph{實作細節}
\label{sec:implementation}
這個有編號的段落示範另一層標題配色。
\subparagraph{更小的細節}
\label{sec:detail}
這一層標題也有編號。
\subsection*{不編號的補充說明}
星號標題維持不編號，也不顯示編譯編號提示。

\section{圖、表與公式}
\label{sec:results}
請參考圖~\ref{fig:flow}、表~\ref{tab:counts} 與公式~\eqref{eq:mean}。
自動引用範例為 \autoref{fig:flow}，文獻引用範例為 \cite{sample2026}。

\begin{figure}[htbp]
  \centering
  \fbox{\rule{0pt}{18mm}\rule{0.65\linewidth}{0pt}}
  \caption{資料處理流程的示意框}
  \label{fig:flow}
\end{figure}
\begin{table}[htbp]
  \centering
  \caption{範例資料筆數}
  \label{tab:counts}
  \begin{tabular}{lr}
    分組 & 筆數 \\
    \hline
    訓練資料 & 80 \\
    驗證資料 & 20 \\
  \end{tabular}
\end{table}
\begin{equation}
  \bar{x} = \frac{1}{n}\sum_{i=1}^{n}x_i
  \label{eq:mean}
\end{equation}
成對的展示型數學分隔符號是合法的：
\[
  \left(\sum_{i=1}^{n}x_i\right)^2 \geq 0.
\]

\section{程式碼與區塊行號}
\label{sec:code}
第一個區塊沿用 style，從第十號開始。各區塊也明確指定語言，方便語法文法辨識。
\begin{lstlisting}[language=Python,firstnumber=10]
def normalize(value):
    cleaned = value.strip()

    return cleaned.lower()
\end{lstlisting}
接續編號，顯示偶數行號，並沿用第一行強制顯示的設定：
\begin{lstlisting}[language=Python,firstnumber=last,stepnumber=2]
values = [" Alpha ", " BETA "]
normalized = [normalize(value) for value in values]
print(normalized)
\end{lstlisting}
重新從一開始，並把行號放在右側：
\begin{lstlisting}[language=Python,numbers=right,firstnumber=1]
for index, value in enumerate(normalized):
    print(index, value)
\end{lstlisting}

\begin{thebibliography}{9}
  \bibitem{sample2026}
  範例作者。\emph{容易閱讀的文件：虛構指南}。2026。
  此為示範資料，並非真實出版品。
\end{thebibliography}
\end{document}
